\documentclass[11pt]{article}

\usepackage[a4paper,margin=1in]{geometry}
\usepackage[T1]{fontenc}
\usepackage{amsmath,amssymb,amsthm}
\usepackage{mathtools}
\usepackage{enumitem}
\usepackage{booktabs}
\usepackage{array}
\usepackage{xcolor}
\usepackage{listings}
\usepackage{algorithm}
\usepackage{algpseudocode}
\usepackage{tikz}
\usetikzlibrary{arrows.meta,positioning}
\usepackage{hyperref}
\usepackage{cleveref}

\hypersetup{colorlinks=true,linkcolor=blue!50!black,citecolor=blue!50!black,urlcolor=blue!50!black}

\newtheorem*{defthm}{Definition}
\newtheorem*{lemthm}{Lemma}
\newtheorem*{thmthm}{Theorem}
\newtheorem*{propthm}{Proposition}
\newtheorem*{remthm}{Remark}

\newcommand{\CP}{\mathbb{CP}}
\newcommand{\type}{\mathrm{type}}
\newcommand{\level}{\mathrm{level}}
\newcommand{\msplit}{\mathrm{split}}
\newcommand{\link}{\mathrm{link}}
\newcommand{\weld}{\mathrm{weld}}
\newcommand{\st}{\mathrm{st}}
\newcommand{\col}{\mathfrak{c}}
\setlength{\emergencystretch}{3em}

\lstdefinestyle{cpp}{
  language=C++,
  basicstyle=\ttfamily\small,
  keywordstyle=\color{blue!60!black}\bfseries,
  commentstyle=\color{green!40!black}\itshape,
  stringstyle=\color{orange!70!black},
  showstringspaces=false,
  breaklines=true,
  columns=flexible,
  tabsize=2,
  frame=single,
  framerule=0.3pt,
  numbers=none,
}

\title{\bf Maubach Subdivision on $\CP^2$}
\author{Vinicius Mello}
\date{\today}

\begin{document}
\maketitle

\begin{abstract}
Maubach's edge-based star-subdivision scheme, as developed in Chapter 4 of
the author's thesis, is a purely combinatorial criterion (the ``Maubach
invariants'') that guarantees a longest-edge-style bisection algorithm
propagates correctly across an entire mesh. The thesis proves this
criterion is satisfied by two specific constructions: the canonical
(Coxeter--Freudenthal--Kuhn) triangulation of a product of simplices, and
the barycentric subdivision of an arbitrary mesh. Both routes rely on a
\emph{global} vertex ranking that is not available in general. This report
verifies the thesis's Lemmas 5--7 (finding and correcting a symbol lost in
the source PDF's text layer along the way), and then shows, with a working
implementation, that the same criterion extends well beyond those two
cases: to meshes in dimension greater than the classical $\le 3$ setting,
and to a triangulation of a closed $4$-manifold with no group-orbit
structure. Its compatibility turns out to be a vertex colouring:
strong compatibility is equivalent to the mesh being balanced. The
application used throughout is the adaptive triangulation of a smooth
algebraic curve $F(X,Y,Z)=0$ -- a Riemann surface -- realized as a
codimension-2 locus inside the complex projective plane $\CP^2$ (a real
4-manifold), triangulated by Gaifullin's 108-cell complex. We describe how
this mesh satisfies Maubach compatibility, how an economical, pointerless
cell representation (Atalay \& Mount's LPT codes) is adapted to a seed with
no algebraic symmetry, and the final version of the resulting adaptive
triangulation algorithm, based on continuation rather than global
refinement.
\end{abstract}

\tableofcontents

% =====================================================================
\section{Introduction}
% =====================================================================

Star-subdivision schemes (Chapter 4 of the thesis~\cite{mello2006})
formalize a simple idea:
a mesh can be refined automatically, guided only by its own combinatorial
data (the partial order on vertices, adjacency, and a hierarchy level),
without ever consulting geometry. Two such schemes are developed there in
parallel -- Mitchell's, based on splitting a distinguished \emph{facet},
and Maubach's, based on splitting a distinguished \emph{edge} (the
classical ``longest-edge'' bisection of Kuhn's canonical triangulation of a
cube). Both are given as an \emph{invariant} $\mathcal I$ (a combinatorial
property a conforming mesh may or may not have) together with a
subdivision algorithm that is proved to preserve $\mathcal I$ forever, once
it holds.

The thesis proves Maubach's invariant is satisfied by exactly two
constructions: the canonical triangulation of a product of $n$ simplices
(Theorem 9 -- of which Kuhn's triangulation of a cube, a product of $n$
intervals, is the special case actually used throughout the rest of the
thesis) and the barycentric subdivision of an arbitrary mesh (Theorem 10).
Both proofs go through the same mechanism: a \emph{global} function that
assigns every vertex of the whole complex a rank in $\{0,\dots,m\}$, such
that every top-dimensional cell sees each rank exactly once. Neither
construction says anything about whether an \emph{arbitrary} triangulation
-- one with no such global ranking built in -- can be made Maubach-compatible.

This report's goal is to show that it can, and that the resulting
machinery genuinely generalizes past the setting the thesis develops it
in: dimension $\le 3$, and meshes derived from Kuhn's cube. We do this
constructively, through a concrete application developed over several
work sessions: the adaptive triangulation of a smooth plane algebraic
curve (a Riemann surface) $F(X,Y,Z)=0$, realized not as a pair of
identified Riemann spheres but directly as a codimension-2 real
submanifold of $\CP^2$ -- a genuinely 4-dimensional real ambient space,
triangulated by A.~Gaifullin's 15-vertex, 108-cell minimal-vertex
triangulation of $\CP^2$~\cite{gaifullin}, a mesh with no product structure
and no flag/rank structure inherited by construction -- although, as
\Cref{sec:gaifullin} shows, it does carry one: it is balanced.

The report is organized as follows. \Cref{sec:maubach} recapitulates the
relevant part of Chapter 4 -- star subdivisions, the Maubach invariants,
and Lemmas 5--7 -- verifying each proof directly (\Cref{sec:maubach}
includes one genuine correction: the PDF's text layer silently drops the
comparison symbols in Lemma 5's statement; \Cref{sec:lemma5} reconstructs
and confirms the correct direction both from the proof and from the
original page image). \Cref{sec:gaifullin} shows how Gaifullin's
triangulation is Maubach-compatible. On a conforming mesh, strong
compatibility is equivalent to the mesh being \emph{balanced}, i.e.\
properly vertex-colourable with $n+1$ colours, and Gaifullin's
triangulation is balanced. Its combinatorics were cross-checked against
GAP's~\cite{gap} \texttt{simpcomp}~\cite{simpcomp} package. \Cref{sec:lpt} describes the
economical, pointerless cell representation this project built on top of
that mesh -- an extension of Atalay \& Mount's LPT
codes~\cite{atalaymount} -- with particular attention to what changes in
the neighbor-finding rules once the seed cells are no longer a single
group orbit. \Cref{sec:continuation} describes, in full detail, the
current (and, as of this writing, final) version of the Riemann-surface
triangulation algorithm, which abandons global refinement in favor of
continuation: following the curve outward from a single confirmed contact
cell, using exactly the neighbor machinery of \Cref{sec:lpt}.

% =====================================================================
\section{Star subdivisions and the Maubach invariant}
\label{sec:maubach}
% =====================================================================

We use the thesis's own notation throughout, restated here only as needed.
A \emph{star subdivision} of a complex $K$ at a simplex $\sigma$, replacing
$\sigma$'s star by a cone from a new vertex $v$ over $\link(\sigma,K)$
joined with $\partial\sigma$, is written $K'=K(\sigma,v)$. An
\emph{ordered} star subdivision $K'=K(\sigma,v)_l$ additionally records
that $v$ occupies position $l$ in every cell of $\st(v,K')$. If $\sigma$ is
itself an $m$-dimensional cell $\langle v_0,\dots,v_m\rangle$ and
$\delta=[j_0,\dots,j_n]\sigma$ is one of its faces, then $K(\delta,v)_l$
produces $n+1$ \emph{child cells} $c_0\sigma,\dots,c_n\sigma$, one per
facet of $\partial\delta$, related by
\[
  \partial_l\, c_i\sigma \;=\; \partial_{j_{n-i}}\,\sigma .
\]
A finite sequence of ordered star subdivisions
$\tau=(\delta_1,v_1)_{l_1}\cdots(\delta_q,v_q)_{l_q}$ imposes a parent/child
hierarchy on cells, tracked by a function $\level(\cdot)$ with
$\level(\sigma)=0$ for $\sigma\in K$ and $\level(c_i\sigma)=\level(\sigma)+1$.

A \emph{subdivision scheme} consists of an invariant $\mathcal I$
(referring only to combinatorial data: the vertex order, adjacency,
$\level$) and an algorithm $\mathrm{Subdivide}(K,\sigma)$ that, to subdivide
a single cell $\sigma$, first propagates a cascade of subdivisions
elsewhere in the mesh as needed to keep $\mathcal I$ true throughout,
before finally splitting $\sigma$ itself. The property used constantly
below is
\begin{equation}
  \partial_i\sigma=\partial_j\omega \;\Longrightarrow\; |\level(\sigma)-\level(\omega)|\le 1 .
  \tag{OneLevel}
\end{equation}

\subsection{The Maubach invariants}

Maubach's scheme~\cite{maubach} splits the edge $[0,\type(\sigma)]$ of a
cell, where
\[
  \type(\sigma) \;:=\; m - \big(\level(\sigma)\bmod m\big) \;\in\;\{1,\dots,m\}.
\]

\begin{defthm}[Definition 9]
An $m$-dimensional mesh $K$ satisfies the \emph{Maubach invariants} if,
besides (OneLevel), for any two adjacent cells $\sigma,\omega$ with
$\partial_i\sigma=\partial_j\omega$,
\begin{align}
  \level(\sigma)=\level(\omega) \;&\Longrightarrow\;
    i=j \quad\text{or}\quad \{i,j\}=\{0,\type(\omega)\}, \tag{MaubachA}\\
  \level(\sigma)=\level(\omega)+1 \;&\Longrightarrow\;
    i=\type(\omega) \ \text{and}\ j\in\{0,\type(\omega)\}. \tag{MaubachB}
\end{align}
\end{defthm}

Writing $\msplit(\sigma):=\sigma[0,\type(\sigma)]$, (MaubachA) says that two
same-level adjacent cells either agree exactly on which vertex their
shared facet omits, or disagree in exactly the pair $\{0,\type(\omega)\}$ --
which, as the thesis shows immediately after the definition, still forces
$\msplit(\sigma)=\msplit(\omega)$. (MaubachB) says that if $\omega$ is one
level shallower than $\sigma$, subdividing $\msplit(\omega)$ will hand one
of $\omega$'s children exactly the facet $\sigma$ already has.

\subsection{Lemma 5: propagation direction, verified}
\label{sec:lemma5}

\begin{lemthm}[Lemma 5]
Let $K$ be a mesh satisfying the Maubach invariants and $\sigma,\omega\in K$
two adjacent cells with $\partial_i\sigma=\partial_j\omega$. Then
\[
\begin{cases}
\level(\omega)\le\level(\sigma), & \text{if } i\notin\{0,\type(\sigma)\}\\
\level(\omega)\ge\level(\sigma), & \text{if } i\in\{0,\type(\sigma)\}.
\end{cases}
\]
Moreover, if $\level(\omega)<\level(\sigma)$, then $i=\type(\omega)$.
\end{lemthm}

\begin{proof}
If $\level(\omega)>\level(\sigma)$, then by (OneLevel)
$\level(\omega)=\level(\sigma)+1$. Applying (MaubachB) with the roles of
the two cells matched so that $\omega$ is the deeper one (i.e.\ with
$\sigma$ playing (MaubachB)'s ``$\omega$''), we get $i\in\{0,\type(\sigma)\}$,
proving the first case (by contraposition: $i\notin\{0,\type(\sigma)\}$
rules out $\level(\omega)>\level(\sigma)$, leaving $\level(\omega)\le\level(\sigma)$
by (OneLevel)). If $\level(\omega)<\level(\sigma)$, then by (OneLevel)
$\level(\sigma)=\level(\omega)+1$, and (MaubachB) applies directly, giving
$i=\type(\omega)$ -- which is the ``moreover'' clause -- and consequently
$i\notin\{0,\type(\sigma)\}$ (since $\type(\cdot)\in\{1,\dots,m\}$ never
takes the value $0$, and $\type(\omega)\ne\type(\sigma)$ because
$\level(\omega)=\level(\sigma)-1$ shifts $\level\bmod m$ by exactly one).
This proves the second case, again by contraposition.
\end{proof}

\begin{remthm}
The PDF text extracted from the thesis's source loses the two comparison
glyphs in the lemma's statement entirely (a font-encoding artifact of the
document, not a content error) -- both lines print as bare
\texttt{level($\omega$)\ \ level($\sigma$)} with the relation missing. We
reconstructed the correct direction independently from the two proof
paragraphs above (the first paragraph derives
$\level(\omega)>\level(\sigma)\Rightarrow i\in\{0,\type(\sigma)\}$, whose
contrapositive gives the \emph{first} printed case; the second derives
$\level(\omega)<\level(\sigma)\Rightarrow i\notin\{0,\type(\sigma)\}$,
giving the \emph{second}), then confirmed it against the original page
image (page 48 of the PDF), which does show $\le$ then $\ge$ exactly as
above. Worth flagging explicitly: this direction is the \emph{opposite} of
the special/generic pairing in Mitchell's analogous Lemma 1 ($i=0\Rightarrow
\level(\omega)\le\level(\sigma)$, $i\ne0\Rightarrow\level(\omega)\ge
\level(\sigma)$ there) -- not a typo, but a real consequence of
(MaubachB)'s built-in asymmetry: unlike Mitchell's fixed special index $0$,
Maubach's special pair always references $\type$ of the \emph{shallower}
cell, regardless of which of $\sigma,\omega$ that is.
\end{remthm}

\subsection{Lemma 6: the child relation}

\begin{lemthm}[Lemma 6]
Let $K$ be a mesh consisting of a single $m$-dimensional cell $\omega$, and
$K'=K(\omega[0,l],v)_l$ with $l>0$. Then the cells $c_0\omega,c_1\omega\in
K'$ satisfy $\partial_0 c_0\omega=\partial_{l-1}c_1\omega$.
\end{lemthm}

\begin{proof}
Write $\omega=\langle v_0,v_1,\dots,v_m\rangle$. Applying $(\omega[0,l],v)_l$
gives
\[
  c_0\omega=\langle v_0,\dots,v_{l-1},v,v_{l+1},\dots,v_m\rangle,\qquad
  c_1\omega=\langle v_1,\dots,v_l,v,v_{l+1},\dots,v_m\rangle,
\]
and dropping position $0$ from $c_0\omega$ or position $l-1$ from
$c_1\omega$ both yield
\[\langle v_1,\dots,v_{l-1},v,v_{l+1},\dots,v_m\rangle.\]
\end{proof}

This is the geometric content of ``splitting the edge $[0,l]$'': the two
children share the facet obtained by dropping the \emph{old} endpoint on
each side, at the position where the \emph{new} vertex $v$ was inserted on
the other side -- exactly the bookkeeping the LPT sibling-facet rule of
\Cref{sec:lpt} will later reuse verbatim, with zero floating-point work.

\subsection{Lemma 7: propagation preserves the invariants}

\begin{lemthm}[Lemma 7]
Let $K$ be a mesh consisting of two adjacent $m$-dimensional cells
$\sigma,\omega$ with $\partial_i\sigma=\partial_j\omega$, such that
$0\le\level(\sigma)-\level(\omega)\le 1$. Then, if $K$ satisfies the
Maubach invariants, so does $K'=K(\omega[0,l],v)_l$ with $l=\type(\omega)$.
\end{lemthm}

\begin{proof}
There are seven possible cases, distinguished by whether $\sigma,\omega$
are at the same level or one apart, and by which of $\{i,j\}$ equal $0$ or
$l$. In every case, $\omega$'s subdivision produces $c_0\omega,c_1\omega$
with $\partial_0c_0\omega=\partial_{l-1}c_1\omega$ (Lemma 6), which already
satisfies (MaubachA) on its own; the seven cases only differ in \emph{which}
additional relation involving $\sigma$ must be checked, and each resolves
to a single facet identity that satisfies either (MaubachA) or (MaubachB)
directly:
\begin{enumerate}[itemsep=1pt,topsep=2pt]
  \item $\level(\sigma)=\level(\omega)$, $i=j=0$: $\partial_0\sigma=\partial_l c_1\omega$, satisfying (MaubachB).
  \item $\level(\sigma)=\level(\omega)$, $i=j=l$: $\partial_l\sigma=\partial_l c_0\omega$, satisfying (MaubachB).
  \item $\level(\sigma)=\level(\omega)$, $i=j=k\notin\{0,l\}$: $\sigma$ is
    also subdivided (into $c_0\sigma,c_1\sigma$, with the same Lemma-6
    relation), and $\partial_k c_0\sigma=\partial_k c_0\omega$,
    $\partial_{k-1}c_1\sigma=\partial_{k-1}c_1\omega$, both satisfying (MaubachA).
  \item $\level(\sigma)=\level(\omega)$, $i=0$, $j=l$: $\partial_0\sigma=\partial_l c_0\omega$, satisfying (MaubachB).
  \item $\level(\sigma)=\level(\omega)$, $i=l$, $j=0$: $\partial_l\sigma=\partial_l c_1\omega$, satisfying (MaubachB).
  \item $\level(\sigma)=\level(\omega)+1$, $i=l$, $j=0$: $\partial_l\sigma=\partial_l c_1\omega$, satisfying (MaubachA).
  \item $\level(\sigma)=\level(\omega)+1$, $i=l$, $j=l$: $\partial_l\sigma=\partial_l c_0\omega$, satisfying (MaubachA).
\end{enumerate}
In every case $K'$ satisfies the Maubach invariants.
\end{proof}

\begin{remthm}
Each of the seven cases was checked directly against Lemma 6's own two
identities and Definition 9's two clauses; none requires anything beyond
substituting the child cell's known facet identity into whichever of
(MaubachA)/(MaubachB) the resulting $\level$/index pair falls under. Cases
1--5 (same level) all end up invoking (MaubachB) precisely because
$\sigma$, unmoved, is now one level \emph{shallower} than $\omega$'s new
children; cases 6--7 invoke (MaubachA) because $\sigma$ was already one
level deeper than $\omega$ and stays so relative to $\omega$'s children.
This is a satisfying cross-check: it is exactly the pattern Lemma 5
predicts for which relation must hold once the levels involved are known.
\end{remthm}

\begin{thmthm}[Theorem 8]
Let $K$ be an $m$-dimensional mesh satisfying the Maubach invariants and
$\sigma\in K$. Then the mesh $K'$ resulting from Maubach's subdivision
algorithm applied to $(K,\sigma)$ also satisfies the Maubach invariants,
and the algorithm is consistent, alternating, and balanced.
\end{thmthm}

The proof is by induction on $\level(\sigma)$, exactly mirroring the
Mitchell case: at the minimal level, Lemma 5 shows every cell sharing
$\sigma$'s subdivision edge $\epsilon=\sigma[0,\type(\sigma)]$ is also at
that level, so the algorithm's own recursive step is never triggered; Lemma
7, applied to every pair of adjacent cells in $\st(\epsilon,K)$, then shows
the resulting mesh still satisfies the invariants. At any higher level,
Lemma 5 and the induction hypothesis reduce to the same base case one level
down. The one detail specific to Maubach (absent from Mitchell's proof) is
that $K$ must already be conforming, since the algorithm implicitly walks
the star of the subdivision edge.

\subsection{Two constructive routes, and their common mechanism}
\label{sec:tworoutes}

The thesis exhibits two families of meshes satisfying Definition 9
outright.

\begin{thmthm}[Theorem 9]
Let $\mathrm{Tr}(\Sigma)$ be the canonical triangulation of a product of
$n$ simplices $\Sigma=\Delta^{m_1}\times\cdots\times\Delta^{m_n}$. Then for
any $\sigma,\omega\in\mathrm{Tr}(\Sigma)$, $\partial_i\sigma=\partial_j\omega
\Rightarrow i=j$; in particular $\mathrm{Tr}(\Sigma)$ satisfies the Maubach
invariants.
\end{thmthm}

\begin{thmthm}[Theorem 10]
Let $K$ be any $n$-dimensional mesh. Then for any $\sigma,\omega\in B(K)$
(its barycentric subdivision), $\partial_i\sigma=\partial_j\omega
\Rightarrow i=j$; in particular $B(K)$ satisfies the Maubach invariants.
\end{thmthm}

Both proofs share one mechanism, worth isolating explicitly because it is
exactly what fails for the mesh this report actually uses. In Theorem 9,
every cell's $j$-th vertex is $v_L$ with $|L|=j$, where $|L|$ is the
\emph{lattice rank} of the multi-index $L\in\{0,\dots,m_1\}\times\cdots
\times\{0,\dots,m_n\}$ -- a quantity defined once, globally, on the whole
vertex set, independent of which cell we look at. In Theorem 10, the
analogous role is played by $d(\delta)$, the dimension of the flag element
$\delta$ labelling a barycentric-subdivision vertex. In both cases, a
shared facet's missing local index is forced to equal the (unique) rank
value absent from that facet's vertex ranks -- so it is automatically the
\emph{same} value on both sides. Call this a \emph{global flag/rank}
structure: a function $\mathrm{rank}:V\to\{0,\dots,m\}$ on the vertex set of
the entire complex such that every top cell's vertices realize every rank
exactly once, with each cell's local vertex order taken to be its rank.
Whenever such a function exists, Maubach compatibility (indeed the
strictly stronger $i=j$ always) follows for free.

Kuhn's own triangulation of a cube is the $n=m$, $m_1=\cdots=m_n=1$ case of
Theorem 9, and it is this special case -- not the general product
construction -- that the rest of the thesis, and the LPT machinery of
\Cref{sec:lpt}, was originally built around. Neither theorem, however, says
anything about a mesh that is \emph{neither} a product of simplices
\emph{nor} a barycentric subdivision -- which is exactly the situation for
the triangulation used in this report. (A related, more elementary
question -- whether the SAME monotone-path mechanism still works when each
\emph{factor} of the product is itself an already-compatible mesh with
more than one cell, rather than a bare simplex -- is answered in
\Cref{app:product}.) Since the thesis was published only in Portuguese,
\Cref{app:proofs} also gives literal English translations of Theorems 8,
9 and 10's own proofs, together with Maubach's subdivision algorithm
itself.

% =====================================================================
\section{Compatibility beyond the cube: balanced seeds}
\label{sec:gaifullin}
% =====================================================================

\begin{remthm}[Revision note]
An earlier version of this section formulated Definition~9 on a seed as
an \emph{edge}-colouring problem on the dual graph. It concluded that
Kühnel's triangulation is Maubach-compatible in the strong sense, and
that Gaifullin's triangulation has no global rank. Both conclusions were
wrong. The dual edge-colouring only requires the omitted vertex to sit at
the same local \emph{index} on both sides of a facet. Definition~9 needs
more: equality $\partial_i\sigma=\partial_j\omega$ of \emph{ordered}
simplices. This section gives the correct formulation, of which the
implementation was an instance all along: the ordering actually used for
Gaifullin's cells is a proper vertex colouring.
\end{remthm}

\subsection{Strong compatibility is balance}

A pure $n$-dimensional complex is \emph{balanced} if its vertices admit a
colouring $\col:V\to\{0,\dots,n\}$ that is injective on every cell.
Equivalently, its $1$-skeleton is properly $(n{+}1)$-colourable.

\begin{propthm}
Let $K$ be a conforming mesh (a hereditary pseudomanifold, as in the
thesis). There is a vertex order in every cell making $K$ strongly
compatible, i.e.\ $\partial_i\sigma=\partial_j\omega\Rightarrow i=j$ as
ordered simplices, if and only if $K$ is balanced. The order is then the
colour order.
\end{propthm}

\begin{proof}
($\Leftarrow$) Order every cell by colour. The shared facet of two
neighbours carries every colour except the one omitted, so $i=j$.
($\Rightarrow$) Equality of ordered facets with $i=j$ puts every shared
vertex at the same position in both cells. The star of a vertex $v$ is
strongly connected, so the position of $v$ is the same in every cell
containing it. Define $\col(v)$ to be this position.
\end{proof}

This is the mechanism of Theorems~9 and~10 (\Cref{sec:tworoutes}) in
general form: the ``global flag/rank'' is precisely a balanced colouring.
Balance is also exactly the classical condition of
Traxler~\cite{traxler1997}, and it is the $N=n$ case of the colouring-based
initialization of Diening, Gehring and Storn~\cite{dgs2025}.

\begin{remthm}[Index-only matching is not enough]
Here is a $3$-dimensional counterexample to the index-only formulation.
Take $\sigma=\langle a,b,c,d\rangle$ and
$\omega=\langle d,x,c,a\rangle$, which share $\{a,c,d\}$ at index $1$ on
both sides, and both have split edge $\{a,d\}$. After one bisection with
new vertex $m$, the children $\langle a,b,c,m\rangle$ and
$\langle x,c,a,m\rangle$ share $\{a,c,m\}$ at indices $1$ and $0$, with
$\type=2$. This violates (MaubachA): the split edge $\{a,c\}$ of the
first lies in the second, whose split edge is $\{x,a\}$. By Kőnig's
theorem, index-only matchings exist whenever the dual graph is
bipartite, so they are far more common than balance.
\end{remthm}

\subsection{Colourability criteria}

Joswig~\cite{joswig2002} characterizes balanced complexes by a trivial
group of projectivities. For simply connected manifolds, this reduces to
Traxler's criterion: every $(n{-}2)$-face has even degree. A cruder
obstruction is often enough: a $2$-neighborly complex with more than
$n+1$ vertices has a complete graph and cannot be balanced. Balance also
implies a chess colouring:

\begin{lemthm}
The dual graph of an orientable, balanced, closed pseudomanifold is
bipartite.
\end{lemthm}

\begin{proof}
Fix a coherent orientation. Let $\varepsilon(\sigma)=\pm1$ record whether
the colour order of $\sigma$ is positively oriented. Two neighbours induce
on their common facet the orientations $(-1)^k\varepsilon(\sigma)$ and
$(-1)^k\varepsilon(\omega)$ relative to its colour order. Coherence
requires these to be opposite.
\end{proof}

\subsection{Gaifullin's triangulation is balanced}

The mesh used throughout this report is A.~Gaifullin's 15-vertex
triangulation of $\CP^2$~\cite{gaifullin}, cross-checked against the
\texttt{simpcomp}~\cite{simpcomp} library of GAP~\cite{gap} (SCLib
\#397). It has $f$-vector $(15,90,240,270,108)$ and Euler
characteristic $3$, and its automorphism group $S_4\times S_3$ acts by
Fubini--Study isometries in Gaifullin's realization.

\begin{propthm}
Gaifullin's triangulation is balanced. Its colour classes are the natural
blocks of its vertex set $(V_4\setminus\{e\})\sqcup(\{1,2,3,4\}\times\{1,2,3\})$:
$V_4\setminus\{e\}$ and $\{a\}\times\{1,2,3\}$ for $a=1,\dots,4$.
Consequently it is a \emph{vertex-minimal balanced} triangulation of
$\CP^2$.
\end{propthm}

\begin{proof}
Every vertex has degree $12$ in the $1$-skeleton, which has $90$ edges,
and its two non-neighbours are the other members of its block. Every
$4$-simplex therefore meets each block exactly once. We checked this
directly on the facet list. Independently, all $240$ triangles have even
degree ($180$ of degree $4$ and $60$ of degree $6$), as Traxler's
criterion requires. For minimality: $\CP^2$ is orientable, so by the
lemma every balanced triangulation of $\CP^2$ admits a chess colouring.
Gaifullin proves that his triangulation is vertex-minimal among
triangulations of $\CP^2$ admitting a chess colouring.
\end{proof}

The local vertex order used by the implementation for the $108$ cells
is this colour order: every vertex occupies the same position in every
cell containing it, and all $270$ shared facets satisfy $i=j$ as ordered
simplices. The bipartiteness of the dual graph, Gaifullin's chess
colouring, which the implementation also uses for coherent orientation,
follows from the lemma. The uniform max/min Fubini--Study edge ratio of
$1.4636$ across all $108$ cells reflects the isometric symmetry.

\subsection{Kühnel's triangulation is not balanced}
\label{sec:kuhnel}

W.~Kühnel and T.~Banchoff's $9$-vertex $\CP^2_9$~\cite{schwartz} is
$3$-neighborly. Its $1$-skeleton is therefore $K_9$, which is not
$5$-colourable, so no vertex order makes it strongly compatible. The
search performed earlier in this project found an index-only matching
(a proper $5$-edge-colouring of its non-bipartite dual graph), which, as
shown above, does not suffice. The barycentric subdivision of $\CP^2_9$
is balanced (Theorem~10). In Kühnel's standard realization, however,
three flag barycentres coincide with existing vertices, because three of
the nine points are sums of two others. This produces degenerate
zero-length-edge cells. Gaifullin's triangulation needs no subdivision
at all.

% =====================================================================
\section{An economical mesh representation: LPT codes and neighbor finding}
\label{sec:lpt}
% =====================================================================

\subsection{Atalay \& Mount's pointerless codes}

A naive implementation of Maubach subdivision keeps every cell ever
created -- bisected or not -- in an append-only container, since other
cells' incidence data may reference it; memory then grows monotonically
with refinement depth. This project's first $\CP^2$ implementation, built
on the library's general-purpose \texttt{nmt<4>} semi-simplicial container,
measured this directly: roughly 4--5\,GB peak resident memory around
subdivision depth 18 on a machine with 7.2\,GB of RAM.

Atalay \& Mount's \emph{LPT code}~\cite{atalaymount} avoids this entirely,
for the specific case of Kuhn's canonical triangulation of a hypercube: a
cell is identified by a small, self-contained integer -- no pointers, no
stored incidence -- from which both the cell's own geometry (their Theorem
1) and its same-\emph{orthant} neighbors' codes (their Theorem 2) are
computable in $O(1)$-ish time by table lookup alone. The key structural
fact their scheme exploits is that Kuhn's $m!$ root simplices are not $m!$
independent objects, but a single orbit of $\mathrm{Sym}(m)$ acting on one
reference simplex; composing a root's signed permutation with the
elementary swap that crosses a same-orthant facet always lands on
\emph{another member of that same orbit}, by group closure -- so a
same-orthant neighbor never needs to consult anything beyond the
permutation-table itself.

\subsection{\texttt{glpt}: extending LPT to a non-orbit seed}

This project's \texttt{glpt} code (\nolinkurl{~/code/lpt/glpt.hpp}, published at \url{https://github.com/vinicius-mello/lpt})
reuses Atalay \& Mount's $(\level,\mathrm{sigperm},\mathrm{orthant\_level})$
machinery verbatim -- it is dimension-generic and independent of which
seed cell a code descends from -- and adds one field, \texttt{seed\_index}
$\in\{0,\dots,107\}$, selecting which of Gaifullin's 108 root cells a given
code belongs to. The resulting cell identity is a handful of bit-packed
fields, all fitting in one 64-bit machine word (\texttt{uint64\_t}):

\begin{center}
\begin{tabular}{@{}lcccccc@{}}
\toprule
field & \texttt{seed} & \texttt{level} & \texttt{sigperm} & \texttt{orthant\_level} & \texttt{orth}[0..3] & total \\
\midrule
width (bits) & 7 & 2 & 9 & 4 & $4\times10=40$ & \textbf{62} \\
range & $0..107$ & $0..3$ & $0..383$ & -- & -- & (of 64) \\
\bottomrule
\end{tabular}
\end{center}

-- 62 of the word's 64 bits, no pointers, no per-cell incidence data
stored anywhere. (\texttt{orth} is not a single field: it is one
$10$-bit coordinate \emph{per dimension} -- $\dim=4$ coordinates, since
one bit is appended to each of them on every subdivision round -- not the
single field a first glance at the table might suggest.) Storage cost is
$O(1)$ per \emph{current leaf}, not per cell ever created
(\Cref{sec:tree} below).

Gaifullin's 108 cells are \textbf{not} a group orbit -- they are 108
independently-specified vertex-order lists, related to each other only by
which facets happen to coincide, a fact established by the search of
\Cref{sec:gaifullin}, not by any algebraic symmetry. So the orbit-closure
property Atalay \& Mount's same-orthant formula relies on simply does not
hold here: empirically, a same-orthant ($1\le i\le\dim$) facet crossing
changes \texttt{seed\_index} in roughly a quarter of sampled cases at
realistic depths, confirmed both statistically and by hand for a concrete
case (a facet whose barycentric weight, relative to the seed's own five
vertices, is exactly zero once the excluded vertex is dropped -- a genuine
boundary touch). A first attempt at reusing the paper's formula unchanged
was cross-checked against the invariant that every cell's vertices must be
convex combinations of its seed's own five original points (since the
whole tree is built by repeated midpointing): it produced a barycentric
weight outside $[0,1]$ -- a geometric impossibility -- in 23.0\% of sampled
same-orthant cases, lining up exactly with the $\sim$24\% of cases that
legitimately do cross a seed boundary.

\subsection{The neighbor function}
\label{sec:neighbor}

\texttt{glpt::neighbor(i,\,\&result)} is entirely combinatorial and uses
no floating point. Its rules are:
\begin{enumerate}
\item \textbf{The sibling facet}: $i=\dim$ for a $0$-child, $i=\type$
  for a $1$-child. Return \texttt{parent().child(other bit)}, by Lemma~6.
\item \textbf{The seed-boundary test}, which takes $O(1)$ time. In
  Atalay \& Mount's frame, the reference root simplex is
  $\{1\ge x_1\ge\dots\ge x_d\ge-1\}$, whose facets lie on $x_1=1$,
  $x_j=x_{j+1}$ and $x_d=-1$. A cell's $k$-th vertex is $(T+y^{(k)})/2^r$.
  Here $y^{(k)}\in\{-1,0,1\}^d$ depends only on (level, signed
  permutation, $k$), and $T_j=2^r-1-2o_j$ depends only on the orthant
  coordinates $o_j$. Facet $i$ lies on root facet $j$ if and only if:
  \begin{itemize}
  \item the linear part of that facet's equation is constant on
    $\{y^{(k)}:k\ne i\}$, a property of (level, signed permutation, $i$)
    alone, stored in a precomputed table ($4\cdot384\cdot5$ entries);
  \item one integer comparison holds: $o_1=0$, or $o_{j+1}-o_j=\delta$,
    or $o_d=2^r-1$.
  \end{itemize}
\item \textbf{Interior facet}: Atalay \& Mount's Theorem~2 formula.
\item \textbf{Seed-crossing facet} on root facet $j$, and every facet of
  a root cell: return the \emph{same code} in the neighbouring seed
  $\mathrm{adj}[s][j]$. This is exact because the seed is balanced. The
  two seeds' barycentric frames agree on their common facet vertex by
  vertex, so the same descent in the other seed contains the same facet,
  at the same local index.
\end{enumerate}
This replaced an earlier version that used two expensive steps: a full
matrix of exact barycentric weights for the boundary test ($O(d^2)$ plus
a loop over the refinement rounds), and a floating-point point location
in the adjacent seed for crossings. The earlier path was the only
source of the two residual disagreements previously reported, at the
maximum representable depth. It also made the
\texttt{GLPT\_DEEP\_CROSSING\_UNIMPLEMENTED} sentinel possible, and that
sentinel can no longer occur.

Verification:
\begin{itemize}
\item the table test agrees with the exact weight test on $7.8$ million
  facets;
\item same-code crossing agrees with geometric location on $1.4$ million
  crossings;
\item \texttt{glpt\_exact\_neighbor\_test} checks $5$ million neighbour
  queries \emph{exactly}, over depths $0$--$43$, using integer barycentric
  weights over the $15$ global vertices. It finds no failure, and no
  non-sibling facet ever changes its local index.
\end{itemize}
Timings: $4.3$\,ns against $79$\,ns per boundary test, and $5$\,ns
against $211$\,ns per crossing. The Riemann-surface program produces
byte-identical output. Orientation coherence at arbitrary depth was
checked on $750\,000$ sampled (cell, facet) pairs, with $0$ violations.

\paragraph{Why codes.}
A pure-topology benchmark (\texttt{mesh\_mem.cpp}) refines Gaifullin's
seed uniformly. It measured $222$ bytes of records per leaf for the
append-only \texttt{nmt<4>} container, which also keeps twice as many
cells as are current, against $19.5$ bytes per leaf for the
\texttt{glpt\_tree} hash table: an $11\times$ reduction before any
per-vertex data is attached.

\subsection{The current-leaf set: \texttt{glpt\_tree}}
\label{sec:tree}

A single packed code identifies a cell, but the \emph{mesh} is the set of
codes that are currently leaves. \texttt{glpt\_tree} stores exactly that
set, as an open-addressing hash table (linear probing, one 64-bit slot per
entry, backward-shift deletion since bisection's continuous
erase-two-insert churn would otherwise accumulate tombstones
indefinitely) -- table occupancy is always exactly the current leaf count,
never a historical operation count, unlike the append-only \texttt{nmt<4>}
container it replaces.

Two operations complete the picture. \texttt{compat\_bisect(r)} bisects
$r$, first recursively forcing whichever neighbors are needed to keep the
mesh \emph{graded} (no facet ever separates cells more than one level
apart) -- a direct, ported ``needs a neighbor finer first'' recursion,
using exactly \texttt{glpt::neighbor()} and \texttt{exists()} the way an
earlier, cube-only implementation of the same idea already did, with three
real bugs found (not merely by inspection but by an exhaustive
graded-invariant test) and fixed along the way: a same-level neighbor that
turns out to be an already-bisected root has no parent to recurse into; a
neighbor whose parent is already mid-recursion higher up the same call
stack must not be processed twice; and the absence of a same-level
neighbor does not always mean it is coarser -- it can instead mean that
cell was already refined \emph{finer}, by the very same top-level call,
moments earlier, through a different facet entirely. Verified after the
fixes at $\sim$40\,000 \texttt{compat\_bisect()} calls cascading to nearly
4 million leaves, 11.6 million facet checks, 0 graded-invariant violations.
\texttt{neighbor\_leaf(r,i,\&result)} then layers the tree's actual state
on top of the bare combinatorial \texttt{neighbor()}: it walks up the
ancestor chain when the combinatorial candidate is not itself a current
leaf (always resolvable, since a cell has exactly one ancestor chain), and
reports \texttt{REFINED\_FURTHER} -- safely skippable, since it means that
region was already visited from elsewhere -- when the candidate has
already been split further than the query.

% =====================================================================
\section{A continuation algorithm for Riemann-surface triangulation}
\label{sec:continuation}
% =====================================================================

\subsection{Problem setup}

A smooth plane curve $F(X,Y,Z)=0$ (homogeneous, degree $n$) sits inside
$\CP^2$ with real codimension 2. Seeding the mesh with Gaifullin's 108
cells and Maubach-bisecting adaptively, a per-cell extraction step
(\texttt{extract\_cell}) turns each visited leaf's boundary root-crossings
-- located by Newton's method on each of its 10 distinct triangular
2-faces, each face solved once and cached by its global vertex-id triple
-- into a graph on those crossing points, decomposed into closed
polygonal cycles that approximate the curve inside that cell. A cell is
called \emph{bad} when that decomposition fails (a vertex of the crossing
graph with degree other than 2, or an unhandled multiple-crossing
configuration on a single facet) -- typically a near-tangency or several
close crossings on one facet.

\subsection{From global refinement to continuation}

The algorithm's earlier refinement criterion, \emph{gradient dispersion}
($1-\kappa$, where $\kappa$ measures how aligned $\nabla F$ is across a
cell's own vertices), is evaluated identically whether or not a cell is
anywhere near the curve, and drives a priority queue seeded from all 108
roots. This refines the entire ambient $\CP^2$ -- a real 4-manifold -- to
resolve a curve of real dimension 2, and is correspondingly wasteful.

The final version of the algorithm instead uses \emph{continuation}:
starting from one cell the curve is confirmed to touch, it explores
outward through only the cells the curve actually reaches, using
precisely the two combinatorial relations of \Cref{sec:lpt} --
\texttt{compat\_bisect}/\texttt{recent\_leaves} to go deeper, and
\texttt{neighbor\_leaf} to move sideways once a cell has reached the
requested depth.

\subsection{The algorithm}

\begin{algorithm}[H]
\caption{Continuation-based Riemann-surface triangulation (default Phase~1
of \texttt{riemann\_cp2\_glpt.cpp})}
\begin{algorithmic}[1]
\Require Gaifullin's 108-cell seed mesh $T$ (all roots present), curve $F$, target depth $D$
\State $s \gets$ first $s\in\{0,\dots,107\}$ with $\mathrm{any\_edge}(\mathrm{extract\_cell}(s))=\mathrm{true}$
\Comment{early-exit scan; skipped entirely for $s'>s$}
\State $\mathit{visited}\gets\varnothing$; $\mathit{frontier}\gets\{s\}$
\While{$\mathit{frontier}\neq\varnothing$}
  \State pop $c$ from $\mathit{frontier}$
  \If{$c\notin T$ \textbf{or} $c\in\mathit{visited}$} \textbf{continue} \Comment{superseded by an earlier cascade}
  \EndIf
  \State $\mathit{visited}\gets\mathit{visited}\cup\{c\}$
  \State $r\gets\mathrm{extract\_cell}(c)$
  \If{\textbf{not} $r.\mathrm{any\_edge}$} \textbf{continue} \Comment{curve doesn't reach here}
  \EndIf
  \If{$\mathrm{simplex\_level}(c) < D$}
    \State $T.\mathrm{compat\_bisect}(c)$
    \State $\mathit{frontier}\gets\mathit{frontier}\cup T.\mathrm{recent\_leaves}()$
    \Comment{includes any cell a conformity cascade touched elsewhere}
  \Else
    \State record $c$ as a final touched leaf (using $r$ already computed)
    \For{$i=0,\dots,\dim$}
      \If{$T.\mathrm{neighbor\_leaf}(c,i)=(\textsc{Found},n)$ \textbf{and} $n\notin\mathit{visited}$}
        \State $\mathit{frontier}\gets\mathit{frontier}\cup\{n\}$
      \EndIf
    \EndFor
  \EndIf
\EndWhile
\end{algorithmic}
\end{algorithm}

The face-crossing cache built while computing $\mathrm{any\_edge}$ above is
shared with the subsequent per-polygon extraction/output pass, so no
2-face crossing is ever solved twice.

\subsubsection*{Correctness sketch}

If the curve crosses cell $A$'s facet $f$, shared with a neighbor $B$, it
must be present, at least infinitesimally, on $B$'s side of $f$ too -- a
1-dimensional curve entering a 4-cell through a shared 3-facet cannot stop
exactly at that facet -- so $\mathrm{any\_edge}(B)$ is guaranteed once the
search reaches it via a genuine crossing facet, and the propagation never
has to guess. This is, however, an argument for recovering a single
\emph{connected component} of $\{F=0\}$: a reducible curve with several
irreducible factors would need one seed per component, which the current
implementation does not attempt.

\subsection{Measured results}

For every one of the six catalog test curves (conic, Fermat cubic, a
generic Hesse-pencil cubic, a smooth quartic, a parabola, and a smooth
elliptic curve -- all connected), continuation recovers the exact same
final leaf set, and the exact same bad-cell count, as an independent
reference that applies the same $\mathrm{any\_edge}$ test to all 108 roots
directly rather than propagating -- checked at depths 8, 12, and 14. Two
further, independent measurements:

\begin{itemize}
\item \textbf{Mesh size}, relative to the earlier gradient-dispersion
  Phase~1 at the same depth and refinement threshold: $62$--$82\%$ of the
  old leaf count at depth 8, falling to $17$--$34\%$ at depth 14 -- a
  $3$--$6\times$ smaller mesh, with the gap widening as depth increases.
\item \textbf{Bad-cell rate}: never worse than the old criterion across
  the six curves, and strictly better on two of them (a conic, $20.13\%
  \to17.81\%$; a parabola, $18.78\%\to17.33\%$) -- exactly the curves
  where the old gradient-dispersion signal had been stopping refinement of
  a touched cell before it reached full depth.
\end{itemize}

A concrete run, the deepest attempted so far: the smooth elliptic curve
$X^2Z-Y^3+YZ^2=0$, seed rotated by a random unitary change of basis (to
avoid the curve resonating with the seed's own coordinate axes), at depth
20, produced 4\,620\,102 leaves, with 222\,981 cleanly-extracted cells and
1\,348 bad ones (a $0.60\%$ bad-cell rate, itself continuing the
already-documented trend that this rate falls with depth). Two additional
targeted repair rounds -- re-bisecting specifically the cells that failed
extraction, up to a small extra depth budget, then re-extracting -- brought
this down to 678 bad cells out of 227\,355 ($0.30\%$).

Re-run at this same depth after \Cref{sec:continuation}'s
\texttt{--tangency-threshold} was added (\Cref{sec:limitations} below),
the full progression for this curve/seed is:

\begin{center}
\begin{tabular}{@{}lrrrr@{}}
\toprule
mode & leaves & ok & bad & bad rate \\
\midrule
baseline & 4\,620\,102 & 222\,981 & 1\,348 & $0.601\%$ \\
\texttt{--repair-rounds 2} alone & 4\,678\,426 & 226\,677 & 678 & $0.298\%$ \\
\texttt{--tangency-threshold 0.05} alone & 6\,788\,844 & 335\,517 & 1\,260 & $0.374\%$ \\
both combined & 6\,842\,524 & 339\,147 & \textbf{610} & $\mathbf{0.180\%}$ \\
\bottomrule
\end{tabular}
\end{center}

-- the same complementary pattern already seen at depth 14, holding at
the deepest depth attempted so far: combining the proactive and reactive
fix roughly halves whichever one alone achieves, for a new low of
$0.180\%$.

The bad-cell rate measures \emph{topological} correctness -- whether a
cell's crossing graph decomposes into simple closed polygons at all -- but
says nothing about how accurately those polygons' own vertices actually
sit on the curve. To check that separately, we evaluated $|F(v)|$ at
every \emph{distinct} extracted crossing point $v$ (deduped by output-
vertex identity, the same identity \Cref{sec:continuation}'s alpha
projection uses for its own global vertex dedup) for this same depth-20
run: 164\,390 distinct points -- matching, as a cross-check, the vertex
count \texttt{--alpha-projection} itself reports for this run -- with
$|F(v)|$ median $1.0\times10^{-16}$, mean $3.0\times10^{-14}$ (pulled up
by a small tail), 95th/99th percentile $2.0\times10^{-13}$/$4.3\times
10^{-13}$, and maximum $9.8\times10^{-13}$. Every one of these values is
within a few orders of magnitude of double-precision machine epsilon
($\approx2.2\times10^{-16}$) relative to $F$'s own natural scale on unit-
norm homogeneous coordinates -- Newton's method (\Cref{sec:continuation}'s
own root-finding, inside \texttt{extract\_cell}) is not merely finding
\emph{a} crossing point per facet, but converging to the true zero locus
to essentially full floating-point precision.

\subsection{Limitations}
\label{sec:limitations}

The algorithm as implemented recovers exactly one connected component of
the curve; a reducible curve needs one seed per irreducible factor, not
yet implemented. A separate investigation, summarized here only for
completeness, ruled out one plausible cause of the residual bad cells: the
geodesic (Fubini--Study) midpoint used for every edge bisection does
\emph{not} split the two children's volumes evenly, unlike the Euclidean
case -- measured directly via a Cayley--Menger volume proxy built from
true geodesic distances -- but that imbalance turns out \emph{not} to
correlate with which cells end up bad; if anything, the bad-cell rate
\emph{falls} as the imbalance of the bisection that created a cell grows.
Bad cells remain concentrated instead on genuine near-tangencies and
multiple close crossings on a single facet -- confirmed quantitatively,
not just asserted, by directly instrumenting \texttt{extract\_cell()}
itself (all six catalog curves, depth 14: 136\,786 touched cells, $8.02\%$
bad):

\begin{itemize}
\item \textbf{Direct cause.} $81.6\%$ of bad cells fail on a
  \emph{degree} mismatch (some crossing-graph node ends up with degree
  $\ne2$, even though every individual facet's own local crossing count
  was already one of $\{0,1,2,4\}$); the remaining $18.4\%$ fail on an
  outright \emph{irregular} facet ($\mathrm{ntn}=3$: three distinct
  crossing points on a single tetrahedral facet's four triangular
  2-faces -- always exactly $3$ in this sample, never more, since
  \texttt{extract\_cell}'s own buffer caps the count at $4$ by
  construction). No cell failed purely at the final cycle-walk step once
  its facet counts and degrees were already consistent.
\item \textbf{Single-facet touching.} A facet with $\mathrm{ntn}=1$ (the
  curve grazing it at one point, without a clean two-point crossing) is
  literally a local near-tangency. Cells with at least one such facet
  have a $29.9\%$ bad rate, versus $4.1\%$ for cells with none -- a
  $\sim7\times$ difference.
\item \textbf{Pairing ambiguity.} The clearest \emph{retrospective}
  signal: for a facet crossed at $\mathrm{ntn}=4$ of its sub-triangles,
  the extraction must choose which $2$ of $3$ candidate pairings connect
  into arcs, by minimum total geodesic length. Taking the ratio of the
  best pairing's cost to the second-best's ($1.0$ = perfectly tied) and
  averaging over every cell with at least one such facet: mean $0.42$
  (median $0.41$) for cells that still decompose cleanly, versus mean
  $0.68$ (median $0.80$, 90th percentile $0.9999$) for bad ones. The bad
  population sits almost exactly where the pairing choice becomes
  ambiguous -- exactly the local signature of a facet close to tangent
  to the curve, where the four crossing points approach a symmetric
  configuration under which every pairing costs nearly the same.
\item \textbf{A direct, pointwise transversality score.} The three
  signals above are all \emph{retrospective}: they need $\mathrm{ntn}\in
  \{1,4\}$ to already have happened. A single crossing root, on its own,
  also carries a direct tangency signal. Write a 2-face's two (chart-
  local, complex) edge vectors as $e_1,e_2$, and the holomorphic gradient
  of $F$ at the root as $(F_a,F_b)$; since $F$ is holomorphic,
  $\nabla\mathrm{Re}(F)$ and $\nabla\mathrm{Im}(F)$ are automatically
  perpendicular and of equal norm (Cauchy--Riemann), so testing whether
  $e_1,e_2$ are close to perpendicular to \emph{both} collapses into one
  clean complex test: with $w_k:=F_a\,e_k^a+F_b\,e_k^b$ ($k=1,2$, the
  ordinary complex directional derivative of $F$ along $e_k$),
  \[
    \mathrm{score} \;=\; \frac{|\mathrm{Im}(\overline{w_1}\,w_2)|}
    {|e_1|\,|e_2|\,(|F_a|^2+|F_b|^2)} \;\in\;[0,1]
  \]
  (Cauchy--Schwarz) is exactly $0$ when the 2-face's own flat tangent
  plane is tangent to the curve at that root (every real direction in it
  has zero directional derivative), and close to $1$ when maximally
  transversal. Taking the minimum over every root a cell examines: bad
  cells skew markedly lower (mean $0.15$, median $0.12$) than good ones
  (mean $0.24$, median $0.21$), and bucketing by score gives an almost
  perfectly monotonic gradient in bad-cell rate: $22.0\%$ for
  score${}\le0.05$, falling to $7.7\%$, $5.6\%$, $3.5\%$, and $2.9\%$
  across the next four buckets up to score $0.8$ (the topmost bucket,
  score${}>0.8$, has only $50$ cells and is not statistically
  meaningful).
\end{itemize}

These four independent checks agree, and together account for the
residual bad-cell population left after continuation and
\texttt{--repair-rounds}.

\paragraph{Acting on it: \texttt{--tangency-threshold}.} Since the
transversality score is computable at a single root, cheaply, from data
\texttt{extract\_cell()} already has on hand, it can be used
\emph{proactively}, unlike the other three (retrospective) signals:
\texttt{riemann\_cp2\_glpt.cpp}'s continuation loop now checks a cell's
own minimum score before accepting it as final at \texttt{--depth}, and
bisects it further -- up to a separate \texttt{--tangency-extra-depth}
budget, mirroring \texttt{--repair-rounds}' extra-depth philosophy but
applied \emph{before} extraction rather than only after it has already
failed. A concrete comparison (elliptic curve, \texttt{--generic}, depth
$14$):

\begin{center}
\begin{tabular}{@{}lrrrr@{}}
\toprule
mode & leaves & ok & bad & bad rate \\
\midrule
baseline & 327{,}370 & 27{,}397 & 887 & $3.14\%$ \\
\texttt{--repair-rounds 2} & 327{,}370 & 30{,}605 & 730 & $2.33\%$ \\
\texttt{--tangency-threshold 0.05} & 612{,}978 & 42{,}278 & 904 & $2.09\%$ \\
both combined & 612{,}978 & 45{,}709 & 729 & $1.57\%$ \\
\bottomrule
\end{tabular}
\end{center}

\texttt{--tangency-threshold} alone cuts the bad-cell rate by about as
much as \texttt{--repair-rounds} alone, at a higher mesh-size cost (it
refines every suspiciously-tangent cell whether or not it would actually
have failed, unlike repair, which only ever touches cells that already
did); combining both -- one proactive, one reactive -- is the best result
measured, roughly halving the baseline rate.

\subsection{Bad-cell rate versus depth}

A sweep over four curves and depths $8$--$16$ (and $18$ for the elliptic
curve) shows a steady decrease, roughly halving every two levels.
Tangency threshold plus repair rounds lowers the rate by a factor of
$1.5$--$3$ at every depth. Curves of higher degree need a few more
levels to reach the same rate. At depth $16$ the combined rate is:
\begin{itemize}
\item $0.46\%$ for the conic;
\item $0.74\%$ for the elliptic curve ($0.39\%$ at depth $18$);
\item $3.1\%$ for the Fermat cubic;
\item $4.3\%$ for the Fermat quartic.
\end{itemize}
The data (\texttt{badrate.csv}) and the plot are in the article's
\texttt{figures/} directory.

\subsection{Topological verification and closing}
\label{sec:topology}

Bad cells leave holes, but the topology can still be verified. The
extracted polygons form a $2$-complex in $\CP^2$, whose vertices are
identified by crossing-node identity: the mesh vertex, edge or $2$-face
the crossing lies on. So the Euler characteristic is available
independently of any projection.

\texttt{--genus-check} proceeds as follows:
\begin{itemize}
\item It splits every \emph{pinched} vertex into one copy per link
  component. A pinched vertex is one whose link has more than one
  component. In all runs, every pinched vertex lies on hole boundaries:
  its link is two paths. About two thirds are a hole touching itself; the
  rest are two holes meeting.
\item It then computes $g=(2c-b-\chi)/2$, where $c$ is the number of
  components and $b$ the number of boundary loops.
\end{itemize}
This recovers the genus $(n-1)(n-2)/2$ exactly for the conic, the
elliptic curve and the Fermat cubic at depth $14$, the elliptic curve at
depth $18$, and the Fermat quartic at depth $17$, even at a bad-cell rate
of $5.9\%$. In particular, every hole is a disc.

\texttt{--close} turns the complex into a closed surface:
\begin{itemize}
\item it splits the pinched vertices;
\item it keeps the main component, dropping islands of $1$ to $18$ faces;
\item it caps every hole with a cone from a \emph{new} apex, placed at
  the Fubini--Study barycentre of the loop and moved onto the curve by
  Newton's method. A fan from an existing vertex could duplicate an
  edge.
\end{itemize}
It then verifies the result: every edge lies in two faces, every link is
one cycle, the surface is connected, and $\chi=2-2g$. The verification
passed in every case above ($\chi=2,0,0,0,-4$), and every cone apex
reached the curve.


% =====================================================================
\section{Conclusion}
% =====================================================================

Maubach's invariants (Definition~9, Chapter~4 of the thesis) are a
purely combinatorial condition. The thesis's two constructive routes,
products of simplices and barycentric subdivision, work by exhibiting a
global vertex rank. This report identifies that rank as a
\emph{balanced colouring}: on a conforming mesh, strong compatibility is
equivalent to balance. The question of which seeds can be used thus
reduces to graph colouring, and it connects with Traxler's condition,
with the colouring-based initialization of Diening, Gehring and
Storn~\cite{dgs2025}, and with the combinatorics of balanced manifolds.
For $\CP^2$, Kühnel's $9$-vertex triangulation is not balanced. Gaifullin's
$15$-vertex triangulation is balanced, and it is a vertex-minimal
balanced triangulation of $\CP^2$.

Balance is also what makes Atalay \& Mount's pointerless LPT codes
portable to arbitrary seeds. A seed index and a seed adjacency table
suffice. The seed-boundary test is a table lookup plus one integer
comparison, and crossing into the next seed keeps the code unchanged.
The resulting mesh uses about $20$ bytes per leaf. On top of it, a
continuation algorithm triangulates smooth plane curves, i.e.\ Riemann
surfaces, directly inside $\CP^2$. Its output is not certified, but its
topology is verified a posteriori, and a topological closing step turns
it into a closed surface of the correct genus.

The material of this report has been reorganized into an article draft
(\nolinkurl{report/article/article.tex}), which also contains a survey of
related work.

% =====================================================================
\appendix
\section{Products of already-compatible meshes}
\label{app:product}
% =====================================================================

This appendix answers a question that arose naturally, not from the $\CP^2$
work of the main text, but from an earlier stage of the same broader
project: the first implementation of \texttt{examples/top/riemann}'s own
Riemann-surface triangulator, which realizes a curve's ambient space not
as $\CP^2$ but as $S^2\times S^2$ (two identified Riemann spheres), seeded
the mesh by triangulating each $S^2$ factor as an \emph{octahedron} (the
boundary of the cross-polytope: 6 vertices, 8 triangular faces) and
combining the two via the classical Eilenberg--Zilber \emph{shuffle}
product -- the same monotone-lattice-path recipe Theorem 9 uses for a
product of simplices, but now with each factor already a small
multi-cell mesh rather than a single simplex. Theorem 9, taken literally,
does not cover this: an octahedron is not a simplex. This raises the
natural question answered below -- does the same construction, applied to
two meshes that are \emph{themselves} already known to satisfy the
thesis's strong compatibility property, still produce a compatible
product?

\begin{propthm}[Product of compatible meshes]
\label{prop:product}
Let $K_1$ be an $m_1$-dimensional mesh equipped with a \emph{global rank}
$\mathrm{rank}_1:V(K_1)\to\{0,\dots,m_1\}$ such that every top cell
$\sigma_1=\langle\sigma_1[0],\dots,\sigma_1[m_1]\rangle$ of $K_1$ satisfies
$\mathrm{rank}_1(\sigma_1[l])=l$ for every $l$ -- i.e., $K_1$ satisfies the
strong form $\partial_i\sigma=\partial_j\omega\Rightarrow i=j$ of the
Maubach invariants, realized (as in Theorems 9 and 10, \Cref{sec:tworoutes})
by an explicit global rank. Let $K_2,\mathrm{rank}_2,m_2$ be likewise. Let
$K_1\times K_2$ be triangulated by the canonical cell-by-cell extension of
the monotone-path construction: for every pair of top cells
$(\sigma_1,\sigma_2)\in K_1\times K_2$, triangulate the prism
$\sigma_1\times\sigma_2$ exactly as in Theorem 9's proof (vertices
$v_L=(\sigma_1[L^1],\sigma_2[L^2])$ for $L=(L^1,L^2)$ in the lattice
$P=\{0,\dots,m_1\}\times\{0,\dots,m_2\}$, cells indexed by monotone paths
in $P$), and glue these triangulated prisms together. Then $K_1\times
K_2$ also satisfies the strong Maubach invariants.
\end{propthm}

\begin{proof}
Define, on the vertex set $V(K_1)\times V(K_2)$ of the product,
\[
  \mathrm{RANK}(a,b) \;:=\; \mathrm{rank}_1(a) + \mathrm{rank}_2(b),
\]
which is well defined globally -- it depends only on the actual vertex
pair $(a,b)$, not on which prism $\sigma_1\times\sigma_2$ we view it
through. Consider any top cell $\tau=\langle v_{L_0},\dots,v_{L_s}\rangle$
of the product triangulation ($s=m_1+m_2$), living in some prism
$\sigma_1\times\sigma_2$, with $L_0,\dots,L_s$ a monotone path in $P$ from
$(0,0)$ to $(m_1,m_2)$. Since each step of a monotone path increases
$L^1+L^2$ by exactly $1$ and $L_0=(0,0)$, we have $|L_l|:=L^1_l+L^2_l=l$
for every $l$. Hence
\[
  \mathrm{RANK}(\tau[l]) \;=\; \mathrm{rank}_1(\sigma_1[L^1_l]) +
  \mathrm{rank}_2(\sigma_2[L^2_l]) \;=\; L^1_l + L^2_l \;=\; l ,
\]
using $\mathrm{rank}_1(\sigma_1[L^1_l])=L^1_l$ and
$\mathrm{rank}_2(\sigma_2[L^2_l])=L^2_l$, which hold by hypothesis on
$K_1,K_2$ individually. So \emph{every} top cell of the glued product --
in any prism whatsoever -- has its $l$-th vertex at global rank exactly
$l$; this did not use anything about $\tau,\tau'$ being in the same or
different prisms.

Now let $\tau,\tau'$ be any two facet-adjacent cells of $K_1\times K_2$,
with $\partial_i\tau=\partial_j\tau'$. The $s$ vertices of this shared
facet have, read off from $\tau$'s own labelling, ranks
$\{0,\dots,s\}\setminus\{i\}$, and read off from $\tau'$'s, ranks
$\{0,\dots,s\}\setminus\{j\}$ -- but $\mathrm{RANK}$ is a single global
function of the vertex, so these are the same set of ranks for the same
set of vertices, forcing $\{0,\dots,s\}\setminus\{i\}=\{0,\dots,s\}
\setminus\{j\}$, hence $i=j$.
\end{proof}

\begin{remthm}
This is exactly Theorem 9's own argument, with $|L|$ replaced by
$\mathrm{RANK}$ -- Theorem 9 is recovered verbatim as the special case
where $K_1,\dots,K_n$ are each a single bare simplex $\Delta^{m_i}$ (whose
``rank'' is simply its own vertex order, with nothing to verify, since a
single simplex has no facet-adjacency at all to check). The genuinely new
content of \Cref{prop:product} is that it also handles $K_1$ or $K_2$
having \emph{more than one cell}, which requires checking that gluing
across different $(\sigma_1,\sigma_2)$ prisms stays consistent -- and the
proof shows this is automatic, precisely because $\mathrm{RANK}$ was
defined once, on the vertex set, rather than separately per prism.
\end{remthm}

\begin{remthm}[Verification on a small example]
Take $K_1$ the path $v_0\!-\!v_1\!-\!v_2$ (two edges), with
$\mathrm{rank}_1(v_0)=\mathrm{rank}_1(v_2)=0$, $\mathrm{rank}_1(v_1)=1$
(forcing the second edge to be oriented $\langle v_2,v_1\rangle$, not
$\langle v_1,v_2\rangle$, to respect this rank), and $K_2$ a single edge
$\langle w_0,w_1\rangle$. This is the smallest instance where one factor
has more than one cell. $K_1\times K_2$ is two unit squares sharing an
edge, each split into two triangles by the usual recipe. Checking all
three non-trivial adjacent pairs by hand -- the two triangles \emph{within}
one square, and the two triangles facing each other \emph{across} the
shared square-to-square edge -- confirms $i=j$ in every case, matching
$\mathrm{RANK}$ exactly as \Cref{prop:product} predicts.
\end{remthm}

\begin{remthm}[Local compatibility is not enough -- a caveat]
\label{rem:localglobal}
\Cref{prop:product}'s hypothesis is that $K_1,K_2$ each admit a
\emph{global} rank, i.e.\ that they are balanced (\Cref{sec:gaifullin}).
Index-only agreement, the same local \emph{index} on both sides of each
facet, is weaker and does not suffice. The $1$-dimensional $3$-cycle
already shows this: pairs of adjacent edges can be ordered to agree, but
the three vertices admit no global $2$-colouring. For Gaifullin's
triangulation the global rank does exist. It is the colouring of
\Cref{sec:gaifullin}.
\end{remthm}

\begin{remthm}[The motivating case: the octahedron]
An octahedron -- vertices $\pm e_1,\pm e_2,\pm e_3$, one face per sign
choice $(\pm,\pm,\pm)$ -- admits an explicit global rank realizing the
strong invariant directly: set $\mathrm{rank}(\pm e_k):=k-1\in\{0,1,2\}$
(color a vertex by \emph{which coordinate axis} it lies on, ignoring
sign). Every face has exactly one vertex on each axis by construction, so
every face is automatically rainbow under this coloring. \Cref{prop:product}
therefore applies directly to $K_1=K_2=$ (octahedron), confirming that the
$S^2\times S^2$ seed mesh's own compatibility -- observed, but not
previously derived from a general statement -- is an instance of a
genuine theorem, not a coincidence particular to that one mesh.
\end{remthm}

% =====================================================================
\section{Translated proofs of Theorems 8--10, and Maubach's algorithm}
\label{app:proofs}
% =====================================================================

The thesis this report is based on~\cite{mello2006} was published only in
Portuguese. For
direct traceability against the original source, this appendix gives
literal English translations of the proofs of Theorems 8, 9 and 10
(stated without proof in \Cref{sec:maubach,sec:tworoutes} of the main
text), together with a full reproduction of Maubach's own subdivision
algorithm (\emph{Listagem 4.2} of the thesis, p.\,47).

\subsection*{Theorem 8}

\begin{thmthm}[Theorem 8]
Let $K$ be an $m$-dimensional mesh satisfying the Maubach invariants and
$\sigma$ a cell of $K$. Then the mesh $K'$ resulting from applying
Algorithm~4.2 to $(K,\sigma)$ also satisfies the Maubach invariants.
Moreover, the algorithm is consistent, alternating, and balanced.
\end{thmthm}

\begin{proof}
By induction on the level of $\sigma$. Let $l$ be the smallest value
assumed by the function $\level$ on $K$. If $\level(\sigma)=l$, it follows
from Lemma 5 that cells adjacent to $\sigma$ that share the subdivision
edge $\epsilon=\sigma[0,\type(\sigma)]$ also have level $l$. Hence the
recursion on line 19 [of Algorithm~4.2] is not executed, and the loop
between lines 8--26 does nothing more than traverse the star of $\epsilon$
and store it in the list $\mathrm{st}$. Applying Lemma 7 to every pair of
adjacent cells $(\sigma',\omega')$, with $\omega'\in\st(\epsilon,K)$, we
get that the subdivision performed on line 28 results in a mesh $K'$
satisfying the statement. If $\level(\sigma)>l$, we can use Lemma 5 and
the induction hypothesis on line 19 to reduce to a case similar to the
initial one.
\end{proof}

\subsection*{Theorem 9}

\begin{thmthm}[Theorem 9]
Let $\mathrm{Tr}(\Sigma)$ be a mesh representing the canonical
triangulation of the product of $n$ simplices $\Sigma=\Delta^{m_1}\times
\Delta^{m_2}\times\cdots\times\Delta^{m_n}$. Then, for any cells
$\sigma,\omega\in\mathrm{Tr}(\Sigma)$,
\[
  \partial_i\sigma=\partial_j\omega \;\Longrightarrow\; i=j.
\]
In particular, $\mathrm{Tr}(\Sigma)$ satisfies the Maubach invariants.
\end{thmthm}

\begin{proof}
A basic property of the construction above is that for every cell
$\sigma\in\mathrm{Tr}(\Sigma)$,
\[
  \sigma[j]=v_L \;\Longleftrightarrow\; |L|=j .
\]
Let $\sigma,\omega\in\mathrm{Tr}(\Sigma)$ be two adjacent cells, with
$\sigma=\langle v_{L_0},\dots,v_{L_s}\rangle$ and $\omega=\langle
v_{L'_0},\dots,v_{L'_s}\rangle$, such that $\partial_i\sigma=\langle
v_{L''_0},\dots,v_{L''_{s-1}}\rangle=\partial_j\omega$. Then there is a
unique $k\in\{0,\dots,s\}$ such that $k\notin\{|L''_0|,\dots,|L''_{s-1}|\}$.
Therefore $|L_i|=k=|L'_j|$, hence $i=k=j$. Note that $k\notin\{0,s\}$,
since for every cell $\sigma\in\mathrm{Tr}(\Sigma)$, $\sigma[0,s]=\langle
v_{(0,\dots,0)},v_{(m_1,\dots,m_n)}\rangle$.
\end{proof}

\subsection*{Theorem 10}

\begin{thmthm}[Theorem 10]
Let $K$ be an $n$-dimensional mesh. Then, for any cells $\sigma,\omega\in
B(K)$,
\[
  \partial_i\sigma=\partial_j\omega \;\Longrightarrow\; i=j.
\]
In particular, $B(K)$ satisfies the Maubach invariants.
\end{thmthm}

\begin{proof}
A basic property of the construction above is that for every cell
$\sigma\in B(K)$,
\[
  \sigma[j]=v_\delta \;\Longleftrightarrow\; d(\delta)=j .
\]
Let $\sigma,\omega\in B(K)$ be two adjacent cells, with $\sigma=\langle
v_{\delta_0},\dots,v_{\delta_n}\rangle$ and $\omega=\langle
v_{\delta'_0},\dots,v_{\delta'_n}\rangle$, such that $\partial_i\sigma=
\langle v_{\delta''_0},\dots,v_{\delta''_{n-1}}\rangle=\partial_j\omega$.
Then there is a unique $k\in\{0,\dots,n\}$ such that $k\notin\{
d(\delta''_0),\dots,d(\delta''_{n-1})\}$. Therefore $d(\delta_i)=k=
d(\delta'_j)$, hence $i=k=j$.
\end{proof}

\subsection*{Maubach's subdivision algorithm}

Below is Maubach's subdivision algorithm (\emph{Listagem 4.2} of the
thesis, p.\,47), given here as pseudocode -- translated in structure and
notation, not transcribed literally from the (mildly archaic,
templated-C++-style) original listing. $k$ is exactly $\type(\sigma)$;
$\mathrm{st}$ (built up as cells are visited) ends up being
$\st(\epsilon,K)$, the star of the subdivision edge
$\epsilon=\sigma[0,k]$ computed at the end; and the recursive call is
precisely the ``propagate to whichever neighbor has a smaller level
first'' step that Theorem 8's proof, above, inducts on.

\begin{algorithm}[H]
\caption{\textsc{MaubachSubdivide}$(K,\sigma)$ (translated from \emph{Listagem 4.2}, p.\,47)}
\begin{algorithmic}[1]
\Require conforming $m$-dimensional mesh $K$, cell $\sigma\in K$ to subdivide
\State $k \gets \type(\sigma)$ \Comment{$=m-(\level(\sigma)\bmod m)$}
\State $\mathrm{st} \gets$ empty list; $\mathit{to\_visit} \gets$ queue containing $\sigma$
\While{$\mathit{to\_visit}$ is not empty}
  \State $c \gets \mathit{to\_visit}.\mathrm{front}()$; $\mathit{to\_visit}.\mathrm{pop}()$
  \If{$c$ is not marked visited}
    \State mark $c$ visited; append $c$ to $\mathrm{st}$
    \For{$i=0,\dots,m$}
      \If{$i=0$ or $i=k$} \textbf{continue} \Comment{these two facets pass to a child whole -- see Lemma 7's own remark}
      \EndIf
      \State $c_a \gets \mathrm{adjacent}(c,i)$
      \If{$c_a$ is not marked visited}
        \If{$\level(c_a) < \level(c)$}
          \State \textsc{MaubachSubdivide}$(K,c_a)$
          \State $c_a \gets \mathrm{adjacent}(c,i)$ \Comment{re-fetch: a child of $c_a$ now occupies this facet}
        \EndIf
        \State push $c_a$ onto $\mathit{to\_visit}$
      \EndIf
    \EndFor
  \EndIf
\EndWhile
\State $\epsilon \gets \sigma[0,k]$ \Comment{the subdivision edge}
\State \textsc{MaubachEdgeSplit}$(K,\epsilon,k,\mathrm{st})$ \Comment{bisects $\epsilon$, splitting every cell of $\mathrm{st}=\st(\epsilon,K)$}
\end{algorithmic}
\end{algorithm}

% =====================================================================
\begin{thebibliography}{99}

\bibitem{mello2006}
Vinícius Moreira Mello.
\emph{Novos Métodos Simpliciais em Computação Gráfica}.
Tese de Doutorado, Instituto de Matemática Pura e Aplicada (IMPA), Rio de
Janeiro, 2006. Advisor: Luiz Velho.

\bibitem{maubach}
J.M. Maubach.
\emph{Local bisection refinement for $N$-simplicial grids generated by
reflection}.
SIAM J. Sci. Stat. Comput., 16:210--227, 1995.

\bibitem{mitchell}
W.F. Mitchell.
\emph{Optimal multilevel iterative methods for adaptive grids}.
SIAM J. Sci. Stat. Comput., 13:146--167, 1992.

\bibitem{gomestavares}
Jonas Gomes and Geovan Tavares.
\emph{Métodos Simpliciais em Computação Gráfica}.
17\textordmasculine\ Colóquio Brasileiro de Matemática, 1989.

\bibitem{lickorish}
W.B.R. Lickorish.
\emph{Simplicial moves on complexes and manifolds}.
In Proceedings of the Kirbyfest, volume 2, pages 299--320, 1999.

\bibitem{atalaymount}
F. Betul Atalay and David M. Mount.
\emph{Pointerless implementation of hierarchical simplicial meshes and
efficient neighbor finding in arbitrary dimensions}.
International Journal of Computational Geometry \& Applications, 2007.

\bibitem{gaifullin}
A.A. Gaifullin.
\emph{A minimal triangulation of complex projective plane admitting a
chess colouring of four-dimensional simplices}.
Proc.\ Steklov Inst.\ Math.\ 266:29--48, 2009. arXiv:0904.4222.

\bibitem{schwartz}
R.E. Schwartz.
\emph{Trisecting the 9-vertex complex projective plane}.
arXiv:2205.00595.

\bibitem{gap}
The GAP Group.
\emph{GAP -- Groups, Algorithms, and Programming}.
\url{https://www.gap-system.org}.

\bibitem{simpcomp}
Felix Effenberger and Jonathan Spreer.
\emph{simpcomp -- a GAP toolbox for simplicial complexes}.
ACM Communications in Computer Algebra, 44(4):186--189, 2010.
arXiv:1004.1367.

\bibitem{traxler1997}
C.T. Traxler.
\emph{An algorithm for adaptive mesh refinement in $n$ dimensions}.
Computing, 59(2):115--137, 1997.

\bibitem{dgs2025}
L. Diening, L. Gehring, J. Storn.
\emph{Adaptive mesh refinement for arbitrary initial triangulations}.
Found.\ Comput.\ Math., 2025. arXiv:2306.02674.

\bibitem{joswig2002}
M. Joswig.
\emph{Projectivities in simplicial complexes and colorings of simple
polytopes}. Math.\ Z., 240(2):243--259, 2002.

\end{thebibliography}

% =====================================================================
\section*{Note on AI assistance}
% =====================================================================

The implementation this report describes -- the \texttt{glpt}/
\texttt{glpt\_tree} pointerless mesh, the Gaifullin-compatibility search,
the neighbor-finding rules of \Cref{sec:lpt}, the continuation algorithm
of \Cref{sec:continuation}, and the various diagnostic investigations
(volume imbalance, bad-cell correlation) referenced along the way -- was
developed with Claude (Anthropic) acting as a coding assistant across many
work sessions, under the author's direction, review, and testing. This
report itself was likewise drafted with Claude's assistance, including
the verification of Lemmas 5--7 in \Cref{sec:maubach}, the translations
and pseudocode of \Cref{app:proofs}, and \Cref{app:product}'s proof --
again under the author's direction and review throughout.

\end{document}
